% =========================================================
% Logical Equivalences
% =========================================================

\section{Logical Equivalences}
% ---------------------------------------------------------

\subsection{Double Negation}

\begin{theorem}[Double Negation]
\[
\neg\neg P \;\equiv\; P
\]
\end{theorem}

\subsection{De Morgan's Laws}

\begin{theorem}[De Morgan's Laws]
\begin{align*}
\neg(P \wedge Q) &\;\equiv\; \neg P \vee \neg Q \\
\neg(P \vee Q) &\;\equiv\; \neg P \wedge \neg Q
\end{align*}
\end{theorem}

\subsection{Commutativity}

\begin{theorem}[Commutativity]
\begin{align*}
P \wedge Q &\;\equiv\; Q \wedge P \\
P \vee Q &\;\equiv\; Q \vee P \\
P \leftrightarrow Q &\;\equiv\; Q \leftrightarrow P
\end{align*}
\end{theorem}

\subsection{Associativity}

\begin{theorem}[Associativity]
\begin{align*}
(P \wedge Q) \wedge R &\;\equiv\; P \wedge (Q \wedge R) \\
(P \vee Q) \vee R &\;\equiv\; P \vee (Q \vee R) \\
(P \leftrightarrow Q) \leftrightarrow R &\;\equiv\; P \leftrightarrow (Q \leftrightarrow R)
\end{align*}
\end{theorem}

\subsection{Distributivity}

\begin{theorem}[Distributivity]
\begin{align*}
P \wedge (Q \vee R) &\;\equiv\; (P \wedge Q) \vee (P \wedge R) \\
P \vee (Q \wedge R) &\;\equiv\; (P \vee Q) \wedge (P \vee R)
\end{align*}
\end{theorem}

\subsection{Idempotence}

\begin{theorem}[Idempotence]
\begin{align*}
P \wedge P &\;\equiv\; P \\
P \vee P &\;\equiv\; P
\end{align*}
\end{theorem}

\subsection{Absorption}

\begin{theorem}[Absorption]
\begin{align*}
P \wedge (P \vee Q) &\;\equiv\; P \\
P \vee (P \wedge Q) &\;\equiv\; P
\end{align*}
\end{theorem}

\subsection{Identity Laws}

\begin{theorem}[Identity Laws]
\begin{align*}
P \wedge \top &\;\equiv\; P \\
P \vee \bot &\;\equiv\; P
\end{align*}
\end{theorem}

\subsection{Domination Laws}

\begin{theorem}[Domination Laws]
\begin{align*}
P \vee \top &\;\equiv\; \top \\
P \wedge \bot &\;\equiv\; \bot
\end{align*}
\end{theorem}

\subsection{Negation Laws}

\begin{theorem}[Negation Laws]
\begin{align*}
P \vee \neg P &\;\equiv\; \top \quad \text{(Law of Excluded Middle)} \\
P \wedge \neg P &\;\equiv\; \bot \quad \text{(Law of Non-Contradiction)}
\end{align*}
\end{theorem}

\subsection{Conditional Equivalences}

\begin{theorem}[Material Implication]
\[
P \rightarrow Q \;\equiv\; \neg P \vee Q
\]
\end{theorem}

\begin{theorem}[Contraposition]
\[
P \rightarrow Q \;\equiv\; \neg Q \rightarrow \neg P
\]
\end{theorem}

\begin{theorem}[Exportation / Importation]
\[
(P \wedge Q) \rightarrow R \;\equiv\; P \rightarrow (Q \rightarrow R)
\]
\end{theorem}

\begin{theorem}[Negation of Conditional]
\[
\neg(P \rightarrow Q) \;\equiv\; P \wedge \neg Q
\]
\end{theorem}

\subsection{Biconditional Equivalences}

\begin{theorem}[Biconditional Expansion]
\[
P \leftrightarrow Q \;\equiv\; (P \rightarrow Q) \wedge (Q \rightarrow P)
\]
\end{theorem}

\begin{theorem}[Biconditional as Disjunction]
\[
P \leftrightarrow Q \;\equiv\; (P \wedge Q) \vee (\neg P \wedge \neg Q)
\]
\end{theorem}

\begin{theorem}[Negation of Biconditional]
\[
\neg(P \leftrightarrow Q) \;\equiv\; P \leftrightarrow \neg Q \;\equiv\; \neg P \leftrightarrow Q
\]
\end{theorem}

\subsection{Duality Principle}

\begin{definition}[Dual Formula]
The \emph{dual} of a formula $\varphi$, denoted $\varphi^d$, is obtained by
simultaneously replacing:
\begin{itemize}
  \item $\wedge$ with $\vee$ and $\vee$ with $\wedge$
  \item $\top$ with $\bot$ and $\bot$ with $\top$
\end{itemize}
while leaving propositional variables and negations unchanged.
\end{definition}

\begin{example}
\begin{align*}
(P \wedge Q)^d &= P \vee Q \\
(P \vee \top)^d &= P \wedge \bot \\
(\neg P \wedge (Q \vee R))^d &= \neg P \vee (Q \wedge R) \\
(P \wedge (P \vee Q))^d &= P \vee (P \wedge Q)
\end{align*}
\end{example}

\begin{theorem}[Duality Principle]
If $\varphi \equiv \psi$ is a logical equivalence involving only the connectives
$\neg$, $\wedge$, $\vee$, $\top$, and $\bot$, then $\varphi^d \equiv \psi^d$ is
also a logical equivalence.
\end{theorem}

\begin{example}[Dual Equivalences]
From each equivalence, we can derive its dual:

\begin{center}
\renewcommand{\arraystretch}{1.3}
\begin{tabular}{|l|l|}
\hline
\textbf{Original Equivalence} & \textbf{Dual Equivalence} \\
\hline
$P \wedge \top \equiv P$ & $P \vee \bot \equiv P$ \\
$P \vee \top \equiv \top$ & $P \wedge \bot \equiv \bot$ \\
$P \wedge (P \vee Q) \equiv P$ & $P \vee (P \wedge Q) \equiv P$ \\
$\neg(P \wedge Q) \equiv \neg P \vee \neg Q$ & $\neg(P \vee Q) \equiv \neg P \wedge \neg Q$ \\
$P \wedge (Q \vee R) \equiv (P \wedge Q) \vee (P \wedge R)$ & $P \vee (Q \wedge R) \equiv (P \vee Q) \wedge (P \vee R)$ \\
\hline
\end{tabular}
\end{center}
\end{example}

\begin{theorem}[Semantic Duality]
For any formula $\varphi$ built from $\neg$, $\wedge$, $\vee$, $\top$, $\bot$:
\[
\varphi^d \equiv \neg\varphi[\neg P_1/P_1, \dots, \neg P_n/P_n]
\]
where $P_1, \dots, P_n$ are the propositional variables in $\varphi$.

Equivalently, $\varphi^d$ evaluated at $(v_1, \dots, v_n)$ equals the negation
of $\varphi$ evaluated at $(\neg v_1, \dots, \neg v_n)$.
\end{theorem}

\subsection{Summary of Logical Equivalences}

\begin{center}
\renewcommand{\arraystretch}{1.3}
\begin{tabular}{|l|l|}
\hline
\textbf{Name} & \textbf{Equivalence} \\
\hline
Double Negation & $\neg\neg P \equiv P$ \\
\hline
De Morgan (1) & $\neg(P \wedge Q) \equiv \neg P \vee \neg Q$ \\
De Morgan (2) & $\neg(P \vee Q) \equiv \neg P \wedge \neg Q$ \\
\hline
Commutativity & $P \wedge Q \equiv Q \wedge P$; \; $P \vee Q \equiv Q \vee P$ \\
\hline
Associativity & $(P \wedge Q) \wedge R \equiv P \wedge (Q \wedge R)$ \\
              & $(P \vee Q) \vee R \equiv P \vee (Q \vee R)$ \\
\hline
Distributivity & $P \wedge (Q \vee R) \equiv (P \wedge Q) \vee (P \wedge R)$ \\
               & $P \vee (Q \wedge R) \equiv (P \vee Q) \wedge (P \vee R)$ \\
\hline
Idempotence & $P \wedge P \equiv P$; \; $P \vee P \equiv P$ \\
\hline
Absorption & $P \wedge (P \vee Q) \equiv P$; \; $P \vee (P \wedge Q) \equiv P$ \\
\hline
Identity & $P \wedge \top \equiv P$; \; $P \vee \bot \equiv P$ \\
\hline
Domination & $P \vee \top \equiv \top$; \; $P \wedge \bot \equiv \bot$ \\
\hline
Negation & $P \vee \neg P \equiv \top$; \; $P \wedge \neg P \equiv \bot$ \\
\hline
Material Implication & $P \rightarrow Q \equiv \neg P \vee Q$ \\
\hline
Contraposition & $P \rightarrow Q \equiv \neg Q \rightarrow \neg P$ \\
\hline
Exportation & $(P \wedge Q) \rightarrow R \equiv P \rightarrow (Q \rightarrow R)$ \\
\hline
Biconditional & $P \leftrightarrow Q \equiv (P \rightarrow Q) \wedge (Q \rightarrow P)$ \\
\hline
\end{tabular}
\end{center}

% ---------------------------------------------------------
